% GENERATED FILE. Edit canonical lessons, then run npm run generate:books.
% canonical-source-hash: fnv1a64:014e0a23accaf214
% canonical-lessons: PT-C56-arvore, PT-C56-flor, PT-C56-folha, PT-C56-montanha, PT-C56-floresta

\chapter{Tree, Flower, Leaf, Mountain, Forest}
\label{ch:pt-tree-flower-leaf-mountain-forest}
\hlchaptermodality{56}

\begin{chapteropening}
\textbf{By the end of this chapter:} \emph{I can name a tree, a flower, a leaf, a mountain and a forest.}

Complete the last lesson of chapter 56: I can name a tree, a flower, a leaf, a mountain and a forest.
\end{chapteropening}

\section[a árvore]{a árvore — a tree}

\label{lesson:PT-C56-arvore}



\begin{quote}
Before the new one: say the Portuguese for to study, then the Portuguese for to travel.
\end{quote}

\subsection*{You'll want to know: a árvore}
\textbf{a árvore} — \textquotedblleft{}a tree\textquotedblright{}.

\begin{cousinweb}
\textbf{árvore} $\leftarrow$ Latin \emph{arbor}, \textquotedblleft{}tree\textquotedblright{} (EtymDB 2.1) — inherited. Old Portuguese said \emph{arvor}.
\end{cousinweb}

\begin{grammarlens}[title={What you've built}]
One more word for this chapter.
\end{grammarlens}

\subsection*{Your turn}
\begin{itemize}
\raggedright
  \item \emph{Say it:} \emph{a árvore}
  \item \emph{Say it:} \emph{a árvore}, once more
  \item \emph{Say it:} say \emph{a árvore}
  \item \emph{Recall:} say the Portuguese for to study, then the Portuguese for to travel, then say \emph{a árvore} again
\end{itemize}

\begin{tcolorbox}[breakable,colback=teal!4,colframe=teal!35!black,title={Before you move on}]
What does a árvore mean? (\textquotedblleft{}A tree\textquotedblright{}.) Say it once more.
\end{tcolorbox}
\section[a flor]{a flor — a flower}

\label{lesson:PT-C56-flor}



\begin{quote}
Before the new one: say the Portuguese for to travel, then the Portuguese for a tree.
\end{quote}

\subsection*{You'll want to know: a flor}
\textbf{a flor} — \textquotedblleft{}a flower\textquotedblright{}.

\begin{cousinweb}
\textbf{flor} $\leftarrow$ Latin \emph{flos, floris}, \textquotedblleft{}flower\textquotedblright{} (EtymDB 2.1) — inherited. It kept its \emph{fl-}, though Latin \emph{fl-} usually became \emph{ch-}.
\end{cousinweb}

\begin{grammarlens}[title={What you've built}]
One more word for this chapter.
\end{grammarlens}

\subsection*{Your turn}
\begin{itemize}
\raggedright
  \item \emph{Say it:} \emph{a flor}
  \item \emph{Say it:} \emph{a flor}, once more
  \item \emph{Say it:} say \emph{a flor}
  \item \emph{Recall:} say the Portuguese for to travel, then the Portuguese for a tree, then say \emph{a flor} again
\end{itemize}

\begin{tcolorbox}[breakable,colback=teal!4,colframe=teal!35!black,title={Before you move on}]
What does a flor mean? (\textquotedblleft{}A flower\textquotedblright{}.) Say it once more.
\end{tcolorbox}
\section[a folha]{a folha — a leaf; a sheet of paper}

\label{lesson:PT-C56-folha}



\begin{quote}
Before the new one: say the Portuguese for a tree, then the Portuguese for a flower.
\end{quote}

\subsection*{You'll want to know: a folha}
\textbf{a folha} — \textquotedblleft{}a leaf; a sheet of paper\textquotedblright{}.

\begin{cousinweb}
\textbf{folha} $\leftarrow$ Latin \emph{folia}, \textquotedblleft{}leaves\textquotedblright{} (EtymDB 2.1), the plural of \emph{folium} — inherited. The Latin plural was heard as a feminine singular.
\end{cousinweb}

\begin{grammarlens}[title={What you've built}]
One more word for this chapter.
\end{grammarlens}

\subsection*{Your turn}
\begin{itemize}
\raggedright
  \item \emph{Say it:} \emph{a folha}
  \item \emph{Say it:} \emph{a folha}, once more
  \item \emph{Say it:} say \emph{a folha}
  \item \emph{Recall:} say the Portuguese for a tree, then the Portuguese for a flower, then say \emph{a folha} again
\end{itemize}

\begin{tcolorbox}[breakable,colback=teal!4,colframe=teal!35!black,title={Before you move on}]
What does a folha mean? (\textquotedblleft{}A leaf; a sheet of paper\textquotedblright{}.) Say it once more.
\end{tcolorbox}
\section[a montanha]{a montanha — a mountain}

\label{lesson:PT-C56-montanha}



\begin{quote}
Before the new one: say the Portuguese for a flower, then the Portuguese for a leaf; a sheet of paper.
\end{quote}

\subsection*{You'll want to know: a montanha}
\textbf{a montanha} — \textquotedblleft{}a mountain\textquotedblright{}.

\begin{cousinweb}
\textbf{montanha} $\leftarrow$ Vulgar Latin \emph{*montanea} (EtymDB 2.1), from Latin \emph{mons, montis}, \textquotedblleft{}mountain\textquotedblright{} — inherited.
\end{cousinweb}

\begin{grammarlens}[title={What you've built}]
One more word for this chapter.
\end{grammarlens}

\subsection*{Your turn}
\begin{itemize}
\raggedright
  \item \emph{Say it:} \emph{a montanha}
  \item \emph{Say it:} \emph{a montanha}, once more
  \item \emph{Say it:} say \emph{a montanha}
  \item \emph{Recall:} say the Portuguese for a flower, then the Portuguese for a leaf; a sheet of paper, then say \emph{a montanha} again
\end{itemize}

\begin{tcolorbox}[breakable,colback=teal!4,colframe=teal!35!black,title={Before you move on}]
What does a montanha mean? (\textquotedblleft{}A mountain\textquotedblright{}.) Say it once more.
\end{tcolorbox}
\section[a floresta]{a floresta — a forest}

\label{lesson:PT-C56-floresta}



\begin{quote}
Before the new one: say the Portuguese for a leaf; a sheet of paper, then the Portuguese for a mountain.
\end{quote}

\subsection*{You'll want to know: a floresta}
\textbf{a floresta} — \textquotedblleft{}a forest\textquotedblright{}.

\begin{cousinweb}
\textbf{floresta} $\leftarrow$ Old French \emph{forest} — borrowed, and reshaped after \emph{flor}. EtymDB 2.1 traces it only to Old Portuguese \emph{foresta}; it derives Old French \emph{forest} from Medieval Latin \emph{forestis}, \textquotedblleft{}open wood\textquotedblright{}.
\end{cousinweb}

\begin{grammarlens}[title={What you've built}]
One more word for this chapter.
\end{grammarlens}

\subsection*{Your turn}
\begin{itemize}
\raggedright
  \item \emph{Say it:} \emph{a floresta}
  \item \emph{Say it:} \emph{a floresta}, once more
  \item \emph{Say it:} say \emph{a floresta}
  \item \emph{Recall:} say the Portuguese for a leaf; a sheet of paper, then the Portuguese for a mountain, then say \emph{a floresta} again
\end{itemize}

\begin{tcolorbox}[breakable,colback=teal!4,colframe=teal!35!black,title={Before you move on}]
What does a floresta mean? (\textquotedblleft{}A forest\textquotedblright{}.) Say it once more.
\end{tcolorbox}
